There is exactly one function from the empty set to itself, namely the empty function, so the counting identity for the number of functions motivates the convention $0^0=1$. However the corresponding two-variable limit is not determined. Along $x=e^{-t},y=\dfrac{1}{t}$ as $t\to\infty$, $x^y=e^{-1}$; along $x=e^{-t},y=\dfrac{1}{t^2}$, it tends to $1$. Thus different paths give different limits. A convention used in combinatorics does not remove this indeterminacy in limits.
